% titel: Blatt 7, A4: Lewis-Strukturen + pi-delokalisierte Bindungen
% Teil II: der durchdachte Weg.

\subsection*{Gegeben und gesucht}
Lewis-Formeln von $\mathrm{NO_3^-}$, $\mathrm{H_2CO}$, $\mathrm{CO_3^{2-}}$,
$\mathrm{O_3}$, $\mathrm{NH_4^+}$, $\mathrm{C_6H_6}$ und Zuordnung, welche
$\pi$-delokalisierte Bindungen haben.

\subsection*{Ansatz}
Lewis-Struktur bauen, Doppelbindungen so verteilen, dass alle Oktette
stimmen; dann Resonanzstrukturen suchen.

\subsection*{Rechnung}
\begin{enumerate}
  \item[a)] \textbf{$\mathrm{NO_3^-}$} (Nitrat, 24 $e^-$): planares N mit
        einer N=O-Doppelbindung und zwei N--O$^-$-Einfachbindungen,
        trigonal-planar. \textbf{3 gleichwertige Resonanzstrukturen},
        Doppelbindung wandernd $\rightarrow$ \textbf{delokalisiert}.
  \item[b)] \textbf{$\mathrm{H_2CO}$} (Formaldehyd): C=O-Doppelbindung,
        C an zwei H und O gebunden (trigonal). Nur \textbf{eine}
        Lewis-Struktur; C=O ist lokalisiert. \textbf{Nicht
        delokalisiert.}
  \item[c)] \textbf{$\mathrm{CO_3^{2-}}$} (Carbonat, 24 $e^-$): C mit
        drei O, eine C=O-Doppelbindung und zwei C--O$^-$; planares
        zentrales C. \textbf{3 Resonanzstrukturen} $\rightarrow$
        \textbf{delokalisiert}.
  \item[d)] \textbf{$\mathrm{O_3}$} (Ozon, 18 $e^-$): gewinkelt
        O--O=O, mit einer O=O-Doppelbindung und einer O--O$-$
        (Formalladung auf End-O). \textbf{2 Resonanzstrukturen}
        $\rightarrow$ \textbf{delokalisiert}.
  \item[e)] \textbf{$\mathrm{NH_4^+}$}: N tetraedrisch an 4 H, alle
        Einfachbindungen, kein freies Paar. Keine
        Doppelbindungen $\rightarrow$ \textbf{nicht delokalisiert}.
  \item[f)] \textbf{Benzol $\mathrm{C_6H_6}$}: hexagonaler Ring aus
        6 C, je eine C=Doppel- und eine C--Einfachbindung abwechselnd.
        \textbf{2 Kekul\'e-Resonanzstrukturen}; das $\pi$-System ist \"uber
        den ganzen Ring delokalisiert $\rightarrow$ \textbf{delokalisiert}.
\end{enumerate}

\begin{ergebnis}
\textbf{$\pi$-delokalisiert:} $\mathrm{NO_3^-}$, $\mathrm{CO_3^{2-}}$,
$\mathrm{O_3}$, Benzol. \textbf{Nicht delokalisiert:}
$\mathrm{H_2CO}$ und $\mathrm{NH_4^+}$.
\end{ergebnis}

\begin{kontrolle}
Resonanz zeigt sich dort, wo eine Doppelbindung an verschiedene
\"aquivalente Positionen wandern kann (Nitrat, Carbonat, Ozon, Benzol).
$\mathrm{H_2CO}$ hat nur eine Stelle f\"ur die C=O, $\mathrm{NH_4^+}$ gar
keine $\pi$-Elektronen.
\end{kontrolle}

\begin{merksatz}
$\pi$-Delokalisierung $\Leftrightarrow$ mesomere Resonanzstrukturen. Wo
eine Doppelbindung um mehrere \"aquivalente Atome ``wandern'' kann, wird
das $\pi$-System delokalisiert und die Bindungen werden gleich lang.
\end{merksatz}
